\subsection{3. Der Text vermischt drei verschiedene
Näherungszustände}\label{der-text-vermischt-drei-verschiedene-nuxe4herungszustuxe4nde}

Die neuen kleinen Matrizen wurden unabhängig diagonalisiert. Sie
bestätigen die gespeicherten numerischen Ritz-Erwartungswerte, aber
nicht deren Umdeutung zu exakten Eigenschaften von Ω.

{\def\LTcaptype{none} % do not increment counter
\begin{longtable}[]{@{}
  >{\raggedright\arraybackslash}p{(\linewidth - 6\tabcolsep) * \real{0.2000}}
  >{\raggedleft\arraybackslash}p{(\linewidth - 6\tabcolsep) * \real{0.2667}}
  >{\raggedleft\arraybackslash}p{(\linewidth - 6\tabcolsep) * \real{0.2667}}
  >{\raggedleft\arraybackslash}p{(\linewidth - 6\tabcolsep) * \real{0.2667}}@{}}
\toprule\noalign{}
\begin{minipage}[b]{\linewidth}\raggedright
Größe
\end{minipage} & \begin{minipage}[b]{\linewidth}\raggedleft
Fünferbasis v0,\ldots,v3,w2
\end{minipage} & \begin{minipage}[b]{\linewidth}\raggedleft
Sechserbasis v0,\ldots,v5
\end{minipage} & \begin{minipage}[b]{\linewidth}\raggedleft
Siebenerbasis zusätzlich w2
\end{minipage} \\
\midrule\noalign{}
\endhead
\bottomrule\noalign{}
\endlastfoot
Ritzenergie/Δ & −1,0942318710 & −1,1384741999 & −1,1384760929 \\
⟨Nb⟩ & 0,8421158543 & 1,0177373262 & 1,0177445059 \\
Gewicht für Addition & 0,02631612045 & 0,03180429144 & 0,03180451581 \\
Gewicht für Entnahme & 0,97368387955 & 0,96819570856 & 0,96819548419 \\
Quadrat der F-Überlappung & hier nicht benötigt & 0,3498924448 &
0,3498903640 \\
Shannon-Entropie der Bosonzahl & hier nicht benötigt & 1,8712372972 Bit
& 1,8712435534 Bit \\
\end{longtable}
}

Die Werte \textbf{0,974/0,026 und 0,03Δ/1,15Δ gehören ausdrücklich zur
kleineren Fünferbasis}, wie \codeword{compute\_charged\_response} ab
Zeile 900 und das Feld \codeword{charged\_response.basis} zeigen. Die
höhere Norm allein macht diese expliziten Vielteilchenvektoren nicht
länger; die Funktion setzt Kc=min(3,Kmax).

Für einen G-Singulettzustand im N=64-Sektor gelten die exakten
Summenregeln

\[\adjustbox{max width=\linewidth}{$\displaystyle 
Z_\mathrm{add}=\langle N_b\rangle/32,\qquad
Z_\mathrm{rem}=1-\langle N_b\rangle/32.
$}\]

Damit sind die verschiedenen Tabellenwerte unmittelbar nachrechenbar.
Die Summe eins ist die CAR-Summenregel für das gesamte spektrale
Gewicht; sie bestimmt keine einzelne Spektrallinie und liefert noch kein
natives Präparations- oder Messinstrument. Der kleine Näherungswert
Zadd≈0,026316 liegt sogar unter dem früheren exakten Ω-Floor
0,842846/32=0,0263389375. Er kann daher nicht als derselbe exakte
Grundzustandswert übernommen werden.

Die Shannon-Entropie der Bosonzahl ist wegen der verschiedenen
Lochzahlen eine Untergrenze für die Boson/Fermion-Verschränkungsentropie
\textbf{dieses jeweiligen reinen Ritz-Zustands}. Der Zahlenwert 1,87 Bit
ist keine hier bewiesene Untergrenze für den wahren nativen
Grundzustand. Das gilt genauso für die 35-Prozent-Überlappung.

\subsubsection{Die mittleren Antwortenergien lassen sich ohne riesige
Zustandsdateien
korrigieren}\label{die-mittleren-antwortenergien-lassen-sich-ohne-riesige-zustandsdateien-korrigieren}

Aus den CAR folgen für jeden normierten G-Singulettzustand mit N=64:

\[\adjustbox{max width=\linewidth}{$\displaystyle 
\begin{aligned}
\sum_r\langle f_r^\dagger Hf_r\rangle
&=\Delta(64\langle N_b\rangle-2\langle N_b^2\rangle)
+2g\operatorname{Re}\langle T_+(62-2N_b)\rangle,\\
\sum_r\langle f_rHf_r^\dagger\rangle
&=2\Delta\langle N_b^2\rangle
+2g\operatorname{Re}\langle T_+(2N_b)\rangle.
\end{aligned}
$}\]

Dabei folgt die erste Identität aus Σf†T+f=T+(Nf−2) und Σf†T−f=T−Nf; die
zweite aus den entsprechenden CAR-Identitäten für f\ldots f†. Die
G-Invarianz macht alle 64 Modenantworten gleich. Nb und die benötigten
T+-Matrixelemente sind bereits in den kleinen Ritzkompressionen
enthalten.

Die separate kleine Rechnung \codeword{late\_a/moment\_audit.py}
reproduziert so die gespeicherten Fünferwerte 0,03107315818Δ und
1,14969200224Δ ohne Millionen Besetzungsamplituden. Für den
\textbf{Siebener-Ritz-Zustand} ergibt dieselbe Rechnung stattdessen die
mittleren Kosten

\[\adjustbox{max width=\linewidth}{$\displaystyle 
\overline\epsilon_{\rm rem}\approx0.03479766990\Delta,\qquad
\overline\epsilon_{\rm add}\approx1.05931330814\Delta.
$}\]

Das korrigiert die Basisverwechslung und ist eine nützliche kleine
Ausleserechnung. Es sind numerische erste Momente eines
Variationszustands, keine neu bewiesenen Ω-Polenergien. Sie dürfen
insbesondere nicht mit der strengen unteren Entnahmeliniengrenze aus
Abschnitt 2 gleichgesetzt werden.

\subsection{4. Was bereits überholt
ist}\label{was-bereits-uxfcberholt-ist}

{\def\LTcaptype{none} % do not increment counter
\begin{longtable}[]{@{}
  >{\raggedright\arraybackslash}p{(\linewidth - 2\tabcolsep) * \real{0.5000}}
  >{\raggedright\arraybackslash}p{(\linewidth - 2\tabcolsep) * \real{0.5000}}@{}}
\toprule\noalign{}
\begin{minipage}[b]{\linewidth}\raggedright
Aussage der Anlage
\end{minipage} & \begin{minipage}[b]{\linewidth}\raggedright
Prüfung
\end{minipage} \\
\midrule\noalign{}
\endhead
\bottomrule\noalign{}
\endlastfoot
Clock sei ein äußerer Handgriff, weil er Gewichte verschiebt & Bereits
in v1.6.6/7 durch das explizite Spin(10)-Wort widerlegt. Nichtskalare
Wirkung innerhalb eines Irreps bedeutet keine äußere Symmetrie. Die
aktuelle Mechanismusprüfung reproduziert zudem den ursprünglichen Clock
und seinen passiven Focklift. \\
Mit allen gewährten inneren Kontrollen bleiben im N=3-Sektor sieben
Kommutantparameter & Bekannte bedingte Algebraaussage; beweist keine
physische Verfügbarkeit dieser Griffe. Die Milliardenangaben sind
Algebradimensionen, keine Zahl physischer Orte. \\
Auf der zweiten Bosonstufe gibt es einen zweiten Zustand & Richtig, aber
unvollständig: v1.6.7 hat bereits genau vier Singulett-Richtungen und
ihre vollständige Projektion berechnet. \\
Der N=64-Grundzustand sei bei g/Δ=1/20 nun wieder grundsätzlich offen &
Die hier benutzten groben Konkurrenzschranken reichen dort nicht; das
widerlegt den früheren stärkeren nativen Satz nicht. Der Quellwert
g*≈0,02705688 gehört zu diesem schwächeren, numerisch ausgewerteten
Schrankenverfahren. \\
Einheitlicher Vermittler müsse universell ein antisymmetrischer Tensor
sein & Der Fünfzyklus ist ein echter bekannter Ausschluss der
angegebenen binären Händigkeit auf primitiven Labels. Die
Feldschlussfolgerung bleibt an diesen ableitungsfreien, unvergrößerten
Ansatz gebunden; kein allgemeines relativistisches Wörterbuch wird
klassifiziert. \\
Alle drei fundamentalen Vorfragen seien jetzt exakt beantwortet & Zu
stark: Kontrollverfügbarkeit, tatsächliche Ω-Erwartungswerte und ein
umfassender relativistischer Feldanschluss bleiben verschieden weit
offen. \\
\end{longtable}
}

\subsection{5. Was daran spannend ist und was als nächstes
trägt}\label{was-daran-spannend-ist-und-was-als-nuxe4chstes-truxe4gt}

Der gute neue Schritt ist \textbf{die tatsächlich abgeschlossene fünfte
Spurnetzwerk-Norm mit der daraus folgenden stärkeren strengen Energie-
und Entnahmeschranke}. Hinzu kommt unsere kleine CAR-Auswertung, die die
geladenen ersten Momente auf den richtigen größeren Ritz-Zustand hebt.
Das verbessert nachprüfbar die Rechnung im bestehenden nativen Modell.

Die Grundlage für die fundamentale Fortsetzung bleibt dagegen der gerade
in v1.6.8 präzisierte gemeinsame Prozess: ursprüngliches
Operationsalphabet, getypter Referenzträger, Zustand, gemeinsame
Auslesung und überprüfbare Interferenz. ν5 ersetzt weder den
Herkunftsnachweis eines zusätzlichen Griffs noch den neuen genauen
Unterschied zwischen gemeinsamer Referenz- dynamik und einem isolierten
Systemantrieb.

\subsection{Reproduktionsbelege}\label{reproduktionsbelege}

\codeword{late\_a/verify.py} prüft elf eingefrorene Eingaben, führt die
ν1-ν5- Spurnetzwerke frisch aus und baut beide Ritzpolynome unabhängig
neu auf. Es gibt \textbf{46 eigene Prüfbedingungen: 36 exakte
beziehungsweise faktische Prüfungen und 10 numerische
Konsistenzprüfungen}, separat dazu 1239 ausgeführte Quellguards der
Normkonstruktion. Die Guardzahl zählt keine unabhängigen mathematischen
Theoreme.

Normal/-OO laufen mit Warnungen als Fehler und liefern byteidentische
Ergebnisdateien mit SHA-256
\codeword{eaa83083d8232d17772d11f8c7c27015a9fe852162586a545d20d39a203f3471}.
Die zusätzliche Momentkontrolle benötigt nur kleine Matrizen und zwei
Quellenpins; sie führt den großen Norm- oder Grundzustandslauf nicht
erneut aus. Kein fremder Lauf wurde als abgeschlossen umetikettiert.

\begin{verbatim}
python3 -B -W error late_a/verify.py --output late_a/results_normal.json
python3 -OO -B -W error late_a/verify.py --output late_a/results_optimized.json
cmp late_a/results_normal.json late_a/results_optimized.json
python3 -B -W error late_a/moment_audit.py --output late_a/moments_normal.json
python3 -OO -B -W error late_a/moment_audit.py --output late_a/moments_optimized.json
cmp late_a/moments_normal.json late_a/moments_optimized.json
\end{verbatim}

\begin{center}\rule{0.5\linewidth}{0.5pt}\end{center}

\section{Nachtragsprüfung B: Feldprojektor und
Lorentz-Händigkeit}\label{nachtragspruxfcfung-b-feldprojektor-und-lorentz-huxe4ndigkeit}

\section{Nachgereichter Text B: Korrekturen und ein neuer prüfbarer
Feldadapter}\label{nachgereichter-text-b-korrekturen-und-ein-neuer-pruxfcfbarer-feldadapter}

Der Text mit 45 Zeilen wurde vollständig gelesen. Seine unveränderte
Kopie steht in \codeword{late\_sources/attachment\_B.txt}, SHA-256
\codeword{9aaf5e18cbc6a84d610b88429872df9b69e6078365bfdc8133c8f17e2c1c3801}.
Die für die neue Prüfung maßgebliche \codeword{field\_dictionary.py} ist
gesondert eingefroren, SHA-256
\codeword{a581bacd17c08dd9a766ae439bfe1e8e35e6d30b9509cb9d98053b897d8ecc13}.
Die bereits abgeschlossenen Schatten-Dateien wurden nicht verändert.

\subsection{Übernahmeentscheidung}\label{uxfcbernahmeentscheidung}

{\def\LTcaptype{none} % do not increment counter
\begin{longtable}[]{@{}
  >{\raggedright\arraybackslash}p{(\linewidth - 2\tabcolsep) * \real{0.5000}}
  >{\raggedright\arraybackslash}p{(\linewidth - 2\tabcolsep) * \real{0.5000}}@{}}
\toprule\noalign{}
\begin{minipage}[b]{\linewidth}\raggedright
Aussage im Text
\end{minipage} & \begin{minipage}[b]{\linewidth}\raggedright
Belastbare Fassung / konkrete Folge
\end{minipage} \\
\midrule\noalign{}
\endhead
\bottomrule\noalign{}
\endlastfoot
1,44 Milliarden unsichtbare Freiheiten im kleinsten nichttrivialen
Sektor & 1444233216 ist die komplexe Kommutantdimension des gewährten
Kontrollsatzes X,Nb im N=3-Sektor. Im N=2-Sektor lautet die
entsprechende Zahl 3829536. Weder Zahl zählt direkt physische
Freiheitsgrade. \\
Mit voller Symmetrie sieben, mit Modenzählern eins; damit nichts
Physisches verborgen & Die assoziativen Algebren sind unter den
ausdrücklich ergänzten Ressourcen bestimmt. Ein skalarer Kommutant
beweist weder vollständige dynamische Lie-Kontrolle noch ausführbare
Präparation/Messinstrumente. Ein neuer exakter Spin-1-Gegenzeuge unten
trennt diese Aussagen. \\
Die Kontrollfrage ist eine einzige Ja/Nein-Frage & Bereits v1.6.6/7
enthalten Zwischenstufen: einzelner Dunkelprojektor, ein getrennter
Casimir, nur SU(4), nur Spin(10). Auch X und Nb unabhängig schalten zu
können ist stärker als ein fixes H. \\
Nur beide getrennten Casimire können die drei dunklen Typen
unterscheiden & Bereits einer der getrennten Casimire genügt, da seine
drei Eigenwerte verschieden sind. Der Gesamtcasimir allein genügt
nicht. \\
Der Skalarkanal ist endgültig verboten; Vermittler muss (1,0) sein &
Exakt im lokalen, ableitungsfreien Ansatz gleicher Weylhändigkeit mit
unverändertem antisymmetrischem W und ohne zusätzliche unabhängige
Komponenten. Der symmetrische Kanal trägt dort. Dies ist kein
vollständiges Feldwörterbuch jedes möglichen Kontinuumsadapters. \\
Das Lochkomposit besitzt einen brauchbaren Spin-1/2-Projektor & Der
angegebene Projektor ist für den \textbf{gleichhändigen}
Tensorproduktraum korrekt. Für das wörtlich adjungierte Weylfeld ist
zunächst die Händigkeit zu klären; die neue exakte Boost-Prüfung unten
macht die Konsequenz sichtbar. \\
N=0 modulo vier verkürzt die Grundzustandskonkurrenz auf
Singulettsektoren & Die Regel gilt für Singuletts. Erst ein unabhängiger
globaler Eindeutigkeitssatz erzwingt für die zusammenhängende
halbeinfache Gruppe einen Singulett-Grundzustand. Ohne ihn kann ein
tieferes nichtsinguläres, entartetes Multiplett gewinnen. Zudem ist die
Liste bosonischer Ladungssektoren nach oben unbeschränkt. \\
480, 916, 299520/229 sind drei exakte Energiesprossen & Es sind
dimensionslose quadrierte Kopplungskoeffizienten, nach Ausklammern von
g², keine drei Eigenenergien. Sie stimmen mit dem Anfang der richtigen
Voll-H-Lanczoskette überein. \\
Nur die gesamte 33-Sprossen-Leiter fehlt noch & 33 zählt die
Bosonzahlstufen k=0,\ldots,32 bei N=64. Eine Zustandsrichtung je Stufe
schließt unter H nicht. Schon k=2 hat vier Singuletts; das
Rücklaufresiduum besitzt exakt positive Norm 5001523200/229. Die
tatsächliche vierte Voll-H-Lanczosdiagonale ist 105168998/26292551·Δ,
nicht 4Δ. \\
Ritzüberlappung 39,7 Prozent bestätigt den wahren Grundzustand & Der
Wert gehört zum benannten Ritz-Zustand. Ohne Residuum und
Spektralabstand ist er keine entsprechende Überlappungsaussage über Ω.
Die zitierte Ritzenergie setzt außerdem g/Δ=1/20 voraus. \\
Schwächere Nachbarsektorfloors lassen Grundzustandssatz wieder offen &
Sie lassen diese schwächere Sonde offen. Das widerlegt oder ersetzt den
stärkeren vorherigen Satz im unveränderten Modellvertrag nicht. \\
Operationssatz und Feldwörterbuch fertig; räumliche Skalierung jetzt
freigegeben & Die erforderlichen Ressourcen und der Lorentzadapter sind
weiterhin konkrete offene Verbindungen. Die neue gemeinsame
Schattenkarte bestimmt deren nächste Prüfung wesentlich genauer als
diese pauschale Freigabe. \\
\end{longtable}
}

Die normierte direkte Summe der drei Feldkanäle hat weiterhin Norm 8√3
bei Normquadrat 192. Die eingefrorene Feldquelle enthält noch die
bereits korrigierte Addition von drei Normen zu 24; ihr gesamtes grünes
Ergebnis wurde daher nicht als unkritischer Vertrauensbeweis übernommen.

\subsection{1. Neuer exakter
Kontrollgegenzeuge}\label{neuer-exakter-kontrollgegenzeuge}

Die gewöhnlichen drei Spin-1-Matrizen J\_x,J\_y,J\_z auf C³ erfüllen
{[}J\_x,J\_y{]}=iJ\_z und zyklisch. Bereits J\_x und J\_z haben
ausschließlich skalare gemeinsame Kommutanten. Der dynamische Lie-Raum
ist trotzdem nur su(2) mit Dimension drei, nicht su(3) mit Dimension
acht. Das wird im neuen Prüfer durch exakte Ränge und alle drei
Kommutatoren kontrolliert.

Damit ist die Unterscheidung praktisch: Ein assoziativer Abschluss kann
vollständig sein, während die tatsächlich durch Hamiltonpulse
erreichbaren Unitären eine echte Untergruppe bilden. Für den nativen
Kontrollvertrag braucht es weiterhin die konkret gewährten Generatoren,
ihre dynamische Lie-Algebra und eine ausgeführte
Puls-/Instrumentenfolge.

\subsection{2. Was am Spin-1/2-Projektor
stimmt}\label{was-am-spin-12-projektor-stimmt}

Die Feldquelle verwendet die sechs unnormierten Koordinaten

\codeword{(00,0), (00,1), (01,0), (01,1), (11,0), (11,1)},

wobei \codeword{(01)} die Summe \codeword{|01>+|10>} bezeichnet. Die
korrekte Gram-Matrix ist

\[\adjustbox{max width=\linewidth}{$\displaystyle 
 G_6=\operatorname{diag}(1,1,2,2,1,1).
$}\]

Der neue Prüfer baut den vollständigen Symmetrisierer auf drei
Spinorindizes unabhängig neu auf. Zugleich konstruiert er die direkte
Epsilon-Auslese

\[\adjustbox{max width=\linewidth}{$\displaystyle 
 C=\begin{pmatrix}0&1&-1&0&0&0\\0&0&0&1&-1&0\end{pmatrix},
 \qquad R=\frac23G_6^{-1}C^T.
$}\]

Exakt gelten

\[\adjustbox{max width=\linewidth}{$\displaystyle 
 CR=I_2,\quad P_{1/2}=RC,\quad
 P_{3/2}=I_6-P_{1/2},\quad
 CP_{3/2}=0,\quad CP_{1/2}=C.
$}\]

Beide Projektoren sind idempotent, komplementär und haben Rang zwei
beziehungsweise vier. P\_\{1/2\} ist selbstadjungiert bezüglich G\_6;
seine nicht symmetrische 6×6-Koordinatenmatrix ist kein Fehler. Die
Epsilon-Auslese und Wiedereinsetzung intertwinen exakt alle drei
sl(2)-Generatoren. Dies liefert eine ausführbare algebraische
Extraktionsregel im \textbf{angenommenen gleichhändigen} Raum

\[\adjustbox{max width=\linewidth}{$\displaystyle 
 (1,0)\otimes(1/2,0)=(3/2,0)\oplus(1/2,0).
$}\]

Die Zahl 1/3 ist nur der Ranganteil im isotropen Zustand. Der neue
Prüfer zeigt positive reine Vektoren mit Spin-1/2-Gewicht null und eins.
Deshalb ersetzt der Projektorrang keine Berechnung seines Gewichts im
nativen Grundzustand oder im geladenen Spektralpol.

